% RESEARCH_DRAFT_SHA256: a25fbdc051dc6d8342f80d906917288b7ee5bb24f5e016d2939103a1a99f8c83
% Layout and source-reference wrappers expanded for the web renderer.
\documentclass[11pt]{article}
\usepackage[margin=25mm]{geometry}
\usepackage{fontspec,unicode-math}
\setmainfont{Latin Modern Roman}
\setmathfont{Latin Modern Math}
\usepackage{amsmath,amssymb,mathtools,xurl,hyperref}
\hypersetup{colorlinks=true,urlcolor=blue,linkcolor=blue}
\setcounter{secnumdepth}{0}
\setlength{\parindent}{0pt}
\setlength{\parskip}{5pt}
\setlength{\emergencystretch}{3em}
\newcommand{\sref}[2]{\hyperlink{source:#1:#2}{[S#2]}}
\newcommand{\cataloguescope}[1]{\paragraph{Source record.}#1}
\newcommand{\floor}[1]{\left\lfloor#1\right\rfloor}
\begin{document}
\section*{Research-attempt provenance}
This standalone research attempt retains the complete original definition body below,
including its historical source and preparation notes. Those notes describe the
earlier catalogue preparation; the new attempt and its verification are identified
separately after that body. No original ID, catalogue status, or source statement is changed.

% UnsolvedMath ID 20001231; source catalogue code AIM-COMPUTATION-0069
\section{Two-sided novelty certificates for sparse binomial output}\label{20001231}
{\small UnsolvedMath ID 20001231 · Computer Science}
\subsection{Definitions and mathematical statement}

\paragraph{Shared notation.}Standard number systems are $\mathbb N,\mathbb Z,\mathbb Q,\mathbb R,\mathbb C$. Specialized notation and index ranges follow the cited source and the definitions below.


A sparse representation stores only nonzero coefficients, monomials, or matrix entries. Algebraic algorithms can produce generating sets that contain redundant elements, and a notion of novelty depends on the mathematical equivalence or invariant being compared. Exact certification of redundancy requires an appropriate membership or equivalence test.
\paragraph{Record scope and qualifications.}
The catalogue title is retained as metadata. Named certificates, contracts, bounds or frameworks in that title are not asserted as results of the original source extract.
The output type, novelty criterion, and desired summary are not specified. The question is retained as an interpretation and software-design prompt, without a particular binomial certificate construction. The supplied catalogue marks this record partially solved; its present research status was not independently checked.
The source text below is rendered literally where necessary; damaged formulas and composite extracts are not silently repaired into a new mathematical claim.
\paragraph{Statement recorded in the supplied source.}
Question 54 (Fienberg). How can we make it easier to identify "new" interesting details among long complicated computer output? Karr: can we take advantage of sparse representations of problems?

Second Open Problem Session \href{https://aimath.org/WWN/compalgstat/compalgstat.pdf}{[S2]}
\paragraph{Source status.}
The supplied catalogue labels this source record partially solved; that metadata was checked against the supplied snapshot on 2026-10-11. The label is preserved for provenance and is not an independent verification of present open status or mathematical correctness.
\subsection{Short English statement}
The source asks how to identify informative new features in large computational output.
\subsection{Sources}
\begin{itemize}
\item[S1] Ulam AI and the UnsolvedMath Contributors, UnsolvedMath: A Curated Collection of Open Mathematics Problems, record 20001231 (AIM-COMPUTATION-0069). Catalogue curation and metadata: CC BY 4.0. The source extract is retained; definitions, English summaries and qualifications are editorial additions.
\url{https://huggingface.co/datasets/ulamai/UnsolvedMath}.
\item[S2] AIM Workshop Problem Lists, source linked by catalogue record AIM-COMPUTATION-0069.
\url{https://aimath.org/WWN/compalgstat/compalgstat.pdf}.
\end{itemize}
\paragraph{Preparation and review.}
The original source and UnsolvedMath contributors retain attribution for the source extract. Codex (GPT-6; exact serving revision not exposed) assisted with definitions, English summaries, scope notes and formatting. No new human mathematical authorship or independently verified solution is asserted. Permanent source ID 20001231 and code AIM-COMPUTATION-0069 checked against the complete supplied registry and current repository. Distinct upstream records, including repeated questions and broad or damaged extracts, are retained for complete coverage of the source CS category. A repeated record is not a new independent conjecture. Complete mathematical-equivalence and current-literature review has not been performed. Preparation and final review follow contribution prompts C and D, with source gaps disclosed. Structural and rendering checks do not establish a new theorem or open status.
\label{end:20001231}

\section{Literature and context check}
AIM Question 54 asks for interpretation of large output and mentions sparse representations, but does not define mathematical novelty. Here novelty is temporarily formalized as nonmembership in a previously known ideal for one elementary family. That restricted criterion is not a definition of scientific interest.
\section{Research attempt: Exact redundancy decisions for a sparse binomial family}
Suppose the known relation in $\mathbb Q[x]$ is $f=x^m-1$ with integer $m\ge1$, and a program outputs $g=x^d-1$ with integer $d\ge1$. Both inputs have two terms, however large their degrees. The exact decision is
\[
g\in(f)\quad\Longleftrightarrow\quad m\mid d.
\]
Write $d=qm+r$, $0\le r<m$. Then
\[
x^d-1=(x^m-1)x^r(1+x^m+\cdots+x^{(q-1)m})+(x^r-1),
\]
where the parenthesized sum is zero if $q=0$. If $r=0$ this is a membership certificate. If $r>0$, the remainder is a nonzero polynomial of degree less than $m$, so cannot be a multiple of the monic degree-$m$ polynomial $f$: a nonzero product $fh$ has degree $m+\deg h$. This proves the converse without numerical root finding.

There is also a witness to failure of set-theoretic implication over $\mathbb C$. Let $\zeta=e^{2\pi i/m}$ when $m>1$. Then $f(\zeta)=0$ but $g(\zeta)=e^{2\pi i r/m}-1\ne0$ for $0<r<m$, since an angle strictly between zero and $2\pi$ does not have exponential one. This is a two-sided audit of this output family: either a quotient identity, or a point contradicting the proposed redundancy. The test uses remainder arithmetic on exponents and need not enumerate all $m$ roots.

Two failures prevent overgeneralizing the audit. First, the expanded quotient certificate can contain $q$ terms although each input has two. For $m=1$ and $d=2^k$, writing every quotient coefficient already requires $2^k$ terms; a compact geometric-sum expression avoids that output cost, but is a different certificate format. Second, equality of zero sets is not general ideal membership. In $\mathbb Q[x]$, the ideal $(x^2)$ and polynomial $x$ have the same zero set, yet $x\notin(x^2)$ by degree. Thus testing roots alone can lose scheme or multiplicity information. The first decision above succeeds because the particular characteristic-zero binomials admit the explicit division identity, not because all sparse outputs are easy to classify.
\section{Verification and adversarial audit}
The division formula is verified by telescoping after multiplication by $x^m-1$. Cases $q=0$, $r=0$, and $m=1$ are treated explicitly. The exponential certificate-size example charges expanded output length rather than claiming an impossible degree-independent printing algorithm.
\paragraph{Audit corrections.} The word novelty is given an explicit provisional ideal-membership meaning; it is not conflated with scientific novelty, support size, or agreement of zero sets.
\paragraph{External computation.} None; symbolic arguments and boundary checks only.
\section{Scope and remaining gaps}
Only redundancy for a selected univariate family has been decided, using known elementary algebra. Multivariate output, statistical significance, and useful scientific novelty remain outside the specified model and unresolved as an agenda.
\paragraph{Integration review.} Codex (GPT-6; exact serving revision not exposed) reviewed the written arguments and expanded layout and source-reference wrappers. Repository metadata and the production web renderer were checked using Python 3.13 and JavaScript with MathJax 3.2.2. This model-assisted review is not independent human proof verification; no PDF compilation or formal proof checking is asserted.

\section{Completion Estimate}
\noindent\textbf{COMPLETION: 0\%}
The calculations below the preserved source statement investigate a concrete route,
but establish no substantiated new progress toward the entire original target.
Known special cases and failed sufficient conditions are not credited as a solution.
\section{References}
The original source references and catalogue attribution are preserved in the
Sources section above. The following URLs were actually consulted for this attempt;
retention of an additional original reference does not claim fresh consultation of it.
\begin{itemize}
\item \url{https://aimath.org/WWN/compalgstat/compalgstat.pdf} (consulted 2026-10-11). 
\end{itemize}
\section{Model and contributor disclosure}
OpenAI gpt-6-astra, requested reasoning configuration ultra, generated this research
attempt on 2026-10-11. The exact serving revision is not exposed. Codex GPT-6
performed orchestration, repository-format review, and a model-assisted mathematical
audit of the written arguments; its exact serving revision
is also not exposed. The mathematical route, boundary checks and final adversarial
self-audit followed contribution prompts A then B in
\texttt{docs/contribution\_prompts.md}, after reading \texttt{CONTRIBUTING.md}.
Python 3.14.2 assembled the standalone documents; any mathematical computation is
specified in Verification. No theorem prover was used. No human mathematical
authorship, substantial human proof correction, or independent human proof checking
is asserted. Model-assisted review, source retrieval, and structural checks do not
establish correctness or novelty. Full private reasoning transcripts are unavailable
as reproducibility artifacts; the complete arguments, inputs and outcomes actually
used are stated here or in the companion batch audit.
\section{Final status}
Unresolved as a new research contribution. This is a concrete route investigation,
not a claim of new complete resolution or a current-status certification. A known
answer reproduced for context does not become a new contribution. The catalogue's
source-status label is preserved; this attempt status does not assert that a
source-labelled solved question is open.
\end{document}
